\documentclass[12pt, a4paper, oneside]{ctexart}
\usepackage{amsmath, amsthm, amssymb, bm, xcolor, tikz, framed, graphicx, hyperref, mathrsfs}
\usepackage{bookmark}
\usepackage{pgfplots}
\pgfplotsset{compat=1.18}


\title{Quantum Mechanics Homework (Week 6)}
\author{Tang Yi \\ [学号已脱敏]}
\usepackage[a4paper, left=2cm, right=2cm, top=2cm, bottom=2cm]{geometry}
\linespread{1.3} % 行距更紧凑
\definecolor{shadecolor}{RGB}{241, 241, 255}

\newcounter{problemname}
\newenvironment{problem}
  {\begin{shaded}\stepcounter{problemname}\par\noindent\textbf{题目 \theproblemname. }}
  {\end{shaded}\par}
\newenvironment{solution}
  {\par\noindent\textbf{题目 \theproblemname 的解答. }}
  {\par}
\newenvironment{note}
  {\par\noindent\textbf{题目 \theproblemname 的注记. }}
  {\par}

\begin{document}

\maketitle

\begin{problem}
  粒子从左边入射, 遇到势场
  \[
  V(x) = 
  \begin{cases}
  -\gamma\delta(x+a), & x<0,\\
  \infty, & x>0.
  \end{cases}
  \]
其中 $a>0$, 求本征波函数。
\end{problem}

\begin{solution}
  设在$\delta$势阱左侧为I区域,右侧非正区域为II区域.在这两区域的波函数演化均满足薛定谔方程 :

  \centerline{$-\dfrac{\hbar^2}{2m} \partial_x^2 \psi = E \psi$}

  因此,在形式上,他们应该有exp的解.

  我们先考虑散射态,即$E > 0$.

  设I区域有解

  \centerline{$\psi_{I} = \mathrm{e}^{ikx} + R\mathrm{e}^{-ikx}$}

  II区域有解,且满足$\psi_{II}(0) = 0$

  \centerline{$\psi_{II} = A \sin(kx)$}

  其中,$k = \dfrac{\sqrt{2mE}}{\hbar}$.

  同时,为了满足连续条件,应该有:
  \begin{align*}
    \psi_{I}(-a) &= \psi_{II}(-a)\\
    \psi_{II}'(-a) - \psi_{I}'(-a) &= -\dfrac{2m\gamma}{\hbar^2} \psi(-a)
  \end{align*}

  代入具体的解,有
  \begin{align*}
    \mathrm{e}^{-ika} + R \mathrm{e}^{ika} &= -A\sin(ka) \\
    Ak\cos(ka) - ik \left(\mathrm{e}^{-ika} - R \mathrm{e}^{ika}\right) &= \dfrac{2m\gamma A}{\hbar^2} \sin(ka)
  \end{align*}

  因此,求解后有
  \begin{align*}
    A &= \dfrac{2ik\mathrm{e}^{-ika}}{k\cos(ka) -(\frac{2m\gamma}{\hbar^2} + ik)\sin(ka)} \\
    R &= -\dfrac{k\cos(ka) +(ik - \frac{2m\gamma}{\hbar^2})\sin(ka)}{k\cos(ka) -(ik + \frac{2m\gamma}{\hbar^2})\sin(ka)} \mathrm{e}^{-2ika}
  \end{align*}

  此时,本征波函数为

  \[
  \psi_{E}(x) = 
  \begin{cases}
    \mathrm{e}^{ikx} + R\mathrm{e}^{-ikx} &,x<-a\\
    A \sin(kx) &,-a<x<0\\
    0 &,x>0
  \end{cases}
  \]

  我们再考虑束缚态,即$E<0$

  设I区域有解

  \centerline{$\psi_{I} = A \mathrm{e}^{-\kappa x}$}

  II区域有解

  \centerline{$\psi_{II} = B \sinh(\kappa x)$}

  其中,$\kappa = \dfrac{\sqrt{-2mE}}{\hbar}$.

  由连续性条件,有
  \begin{align*}
    A \mathrm{e}^{\kappa a} &= -B \sinh{\kappa a}\\
    \kappa A \mathrm{e}^{\kappa a} + \kappa B \cosh(\kappa a) &= \dfrac{2m\gamma B}{\hbar^2} \sinh(\kappa a)
  \end{align*}

  化简后,有

  \centerline{$\dfrac{2m\gamma}{\hbar^2} = \kappa (\coth(\kappa a) -1)$}

  若该方程有解$\kappa_n$,则说明存在束缚态能级$E_n = -\dfrac{\hbar^2 \kappa_n^2}{2m}$.反之,则不存在.

  此时本征波函数为

  \[
  \psi_n(x)=
  \begin{cases}
    A \mathrm{e}^{\kappa_n x} &,x<-a\\
    -A \dfrac{\mathrm{e}^{\kappa_n a}}{\sinh(\kappa_n a)} \sinh(\kappa_n x) &,-a<x<0\\
    0 &,x>0
  \end{cases}
  \]

  其中,$A$由归一化条件决定.

  \centerline{$A = \sqrt{\dfrac{2\kappa_n}{1 - \mathrm{e}^{-2\kappa_n a} + \kappa_n a \coth(\kappa_n a) - \kappa_n a}}$}
\end{solution}
  
\begin{problem}
  证明：若势能函数为复，即 $V(r) = V_1(r) + iV_2(r)$ ，则薛定谔方程无法满足几率守恒条件。
\end{problem}

\begin{solution}
  我们如果还采用哥本哈根诠释,把波函数模长平方解释为概率密度,即有

  \centerline{$\rho = \psi^* \psi$}

  则
  \begin{align*}
    \partial_t \rho &= \partial_t \psi^* \psi + \psi^* \partial_t \psi \\
    &= -\dfrac{H^*}{i\hbar} \psi^* \psi +\psi^* \dfrac{H}{i\hbar} \psi \\
    &= \dfrac{i\hbar}{2m} (\psi^* \nabla^2 \psi - \nabla^2 \psi^* \psi) +\dfrac{1}{i\hbar} (\psi^* V \psi - V^* \psi^* \psi) \\
    &= -\nabla \cdot \bm{J} + \dfrac{2V_2}{\hbar}  \psi^* \psi
  \end{align*}

  显然,由于虚势的存在,流守恒方程被打破.薛定谔方程无法满足几率守恒条件.
\end{solution}

\begin{problem}
  尽管波函数中的普乘相因子常量没有任何物理意义，但当考虑态叠加时，相对相因子却是有意义的。考虑无限深方势阱，阱底范围$0\sim a$，若初始时粒子处于状态：

  \centerline{$\Psi(x,0) = A[\psi_1(x)+\mathrm{e}^{i\phi}\psi_2(x)]$}
其中$A,\phi$是常数， $\psi_n(x)$ 是归一化的本征态。求 $\Psi(x,t)$,$|\Psi(x,t)|^2$ 和 $\langle x\rangle$。分别
讨论 $\phi= 0$，$\phi= \dfrac{\pi}{2}$ 和 $\phi= \pi$ 的具体情况。
\end{problem}

\begin{solution}
  由归一化条件,不难得到$A = \dfrac{1}{\sqrt{2}}$.加上时间因子,不难有

  \centerline{$\Psi(x,t) = \dfrac{1}{\sqrt{2}} [\psi_1(x)\mathrm{e}^{-iE_1t/\hbar} + \psi_2(x)\mathrm{e}^{i(-E_2t/\hbar+\phi)}]$}

  对其求模方,有

  \centerline{$\Psi^* \Psi = \dfrac{1}{2} \left[\psi_1^* \psi_1 + \psi_2^*\psi_2 + \psi_1^*\psi_2 \mathrm{e}^{i((E_1 - E_2)t/\hbar+\phi)} + \psi_2^*\psi_1 \mathrm{e}^{-i((E_1-E_2)t/\hbar + \phi)}\right]$}

  由于是无限深势阱,所以本征波函数是实的.整理后

  \centerline{$|\Psi(x,t)|^2 = \dfrac{1}{2}\left[\psi_1^2 + \psi_2^2 + 2\psi_1\psi_2 \cos\left(\dfrac{(E_1 - E_2)t}{\hbar}+\phi\right)\right]$}

  其中,$E_n = \dfrac{n^2\pi^2\hbar^2}{2ma^2}$,$\psi_n = \sqrt{\dfrac{2}{a}}\sin\dfrac{n\pi x}{a}$.

  因此,代入后为

  \centerline{$|\Psi(x,t)|^2 = \dfrac{1}{2}\left[\psi_1^2 + \psi_2^2 + 2\psi_1\psi_2 \cos\left(\dfrac{3\pi^2\hbar}{2ma^2}t-\phi\right)\right]$}

  然后求$\langle x \rangle$,
  \begin{align*}
    \langle x \rangle &= \langle \Psi| x | \Psi \rangle\\
    &= \dfrac{1}{2} \left\langle \left.\psi_1(x)\mathrm{e}^{-iE_1t/\hbar} + \psi_2(x)\mathrm{e}^{i(-E_2t/\hbar+\phi)}\right. \right| x \left| \left.\psi_1(x)\mathrm{e}^{-iE_1t/\hbar} + \psi_2(x)\mathrm{e}^{i(-E_2t/\hbar+\phi)} \right. \right\rangle \\
    &= \dfrac{1}{2} \left( \dfrac{a}{2} + \dfrac{a}{2} + \mathrm{e}^{i[(E_1 - E_2)t/\hbar + \phi]} \langle \psi_1 | x | \psi_2 \rangle + \mathrm{e}^{-i[(E_1 - E_2)t/\hbar+\phi]} \langle \psi_2 | x | \psi_1 \rangle\right) \\
    &= \dfrac{a}{2} + \langle \psi_1 | x | \psi_2 \rangle \cos\left( \dfrac{3\pi^2\hbar}{2ma^2}t-\phi\right) \\ 
    &= \dfrac{a}{2} - \dfrac{16a}{9\pi^2}\cos\left( \dfrac{3\pi^2\hbar}{2ma^2}t-\phi\right) 
  \end{align*}

  可以看到$\phi$的不同实际上导致干涉项相位的初值不同.结果如下所示.

  \[
\begin{array}{c|c|c|c}
\hline
\phi & \Psi(x,t) & |\Psi(x,t)|^2 & \langle x\rangle \\[2pt]
\hline
0 &
\dfrac{1}{\sqrt2}\big(\psi_1 e^{-iE_1t/\hbar}+\psi_2 e^{-iE_2t/\hbar}\big) &
\dfrac{1}{2}(\psi_1^2+\psi_2^2)+\psi_1\psi_2\cos(\frac{3\pi^2\hbar}{2ma^2} t) &
\dfrac{a}{2}-\dfrac{16a}{9\pi^2}\cos(\frac{3\pi^2\hbar}{2ma^2} t)
\\[8pt]
\dfrac{\pi}{2} &
\dfrac{1}{\sqrt2}\big(\psi_1 e^{-iE_1t/\hbar}+i\psi_2 e^{-iE_2t/\hbar}\big) &
\dfrac{1}{2}(\psi_1^2+\psi_2^2)-\psi_1\psi_2\sin(\frac{3\pi^2\hbar}{2ma^2} t) &
\dfrac{a}{2}-\dfrac{16a}{9\pi^2}\sin(\frac{3\pi^2\hbar}{2ma^2} t)
\\[8pt]
\pi &
\dfrac{1}{\sqrt2}\big(\psi_1 e^{-iE_1t/\hbar}-\psi_2 e^{-iE_2t/\hbar}\big) &
\dfrac{1}{2}(\psi_1^2+\psi_2^2)-\psi_1\psi_2\cos(\frac{3\pi^2\hbar}{2ma^2} t) &
\dfrac{a}{2}+\dfrac{16a}{9\pi^2}\cos(\frac{3\pi^2\hbar}{2ma^2} t)
\\[4pt]
\hline
\end{array}
\]

\end{solution}

\begin{problem}
  利用基本对易关系和对易关系恒等式，计算对易式:1)$[\hat{r}_x,\hat{T}]$;2)$[\hat{l}_x,\hat{l}_y]$;3)$[\hat{l}^2,\hat{l}_y]$;\\
  4)$[\hat{l}_+,\hat{l}_-]$.
\end{problem}

\begin{solution}
    为了方便,答案中的力学量算符都略去hat.

  (1)
  \begin{align*}
    [r_x,T] &= \left[r_x,\dfrac{p_x^2+p_y^2+p_z^2}{2m}\right] \\
    &= \dfrac{1}{2m} [r_x,p_x^2] \\
    &=\dfrac{1}{2m}(p_x[r_x,p_x]+[r_x,p_x]p_x)\\
    &=\dfrac{i\hbar}{m}p_x
  \end{align*}

  (2)
  \begin{align*}
    [l_x,l_y] &= [\varepsilon_{xij}r^ip^j,\varepsilon_{ykl}r^kp^l]\\
    &=\varepsilon_{xij}\varepsilon_{ykl}[r^ip^j,r^kp^l]\\
    &=\varepsilon_{xij}\varepsilon_{ykl}([r^ip^j,r^k]p^l+r^k[r^ip^j,p^l])\\
    &=\varepsilon_{xij}\varepsilon_{ykl}(r^i[p^j,r^k]p^l+r^k[r^i,p^l]p^j)\\
    &=\varepsilon_{xij}\varepsilon_{ykl}(-i\hbar\delta^{jk}r^ip^l+i\hbar\delta^{il}r^kp^j)\\
    &=i\hbar(\varepsilon_{xij}\varepsilon_{yk}\,^{i}r^kp^j - \varepsilon_{xij}\varepsilon_{y}\,^j \,_{l}r^ip^l)\\
    &=i\hbar[(\delta_{jy}\delta_{xk}-\delta_{jk}\delta_{xy})r^kp^j-(\delta_{xl}\delta_{iy}-\delta_{xy}\delta_{il})r^ip^l]\\
    &=i\hbar(\delta_{xi}\delta_{yl} - \delta_{xl}\delta_{yi})r^ip^l\\
    &=i\hbar \varepsilon^z\,_{xy} \varepsilon_{zil}r^ip^l\\
    &=i\hbar \varepsilon^z\,_{xy}l_z = i\hbar l_z
  \end{align*}

  (3)
  \begin{align*}
    [l^2,l_y] &= [l^il_i,l_y]\\
    &=l^i[l_i,l_y]+[l^i,l_y]l_i\\
    &=l^i i\hbar \varepsilon_{iyk} l^k + i\hbar \varepsilon^i\,_{yk} l^k l_i \\
    &=i\hbar \varepsilon_{iyk}(l^il^k + l^kl^i)\\
    &=2i\hbar \varepsilon_{y[ki]}l^{(ik)} = 0
  \end{align*}

  (4)
  \begin{align*}
    [l_+,l_-] &= [l_x + il_y,l_x - il_y] \\
    &=[l_x,-il_y]+[il_y,l_x]\\
    &=-2i [l_x,l_y]\\
    &=-2i (i\hbar l_z) = 2\hbar l_z
  \end{align*}
\end{solution}
\end{document}